\fancyhead[L]{}
\fancyhead[R]{\textsf{\textbf{\color{stewartcyan}CHƯƠNG 1}}\quad Ôn tập\quad \textbf{69}}

\noindent
\begin{minipage}[t]{0.485\textwidth}
\begin{enumerate}[label=\textbf{\arabic*.}, leftmargin=*, start=24]
    \item Một nhà sản xuất thiết bị gia dụng nhỏ nhận thấy rằng chi phí để sản xuất $1000$ lò nướng bánh mì mỗi tuần là $\$9000$ và chi phí để sản xuất $1500$ lò nướng bánh mì mỗi tuần là $\$12\,000$.
    \begin{enumerate}[label=(\alph*), leftmargin=*]
        \item Biểu diễn chi phí như một hàm số theo số lượng lò nướng được sản xuất, giả sử rằng đó là một hàm tuyến tính. Sau đó vẽ đồ thị của hàm số.
        \item Hệ số góc của đồ thị là bao nhiêu và nó đại diện cho điều gì?
        \item Giao điểm của đồ thị với trục tung ($y$-intercept) là bao nhiêu và nó đại diện cho điều gì?
    \end{enumerate}

    \item Cho $f(x) = 2x + 4^x$, hãy tìm $f^{-1}(6)$.

    \item Tìm hàm số ngược của $f(x) = \dfrac{2x + 3}{1 - 5x}$.

    \item Sử dụng các quy tắc của logarit để khai triển mỗi biểu thức sau.
    \begin{enumerate}[label=(\alph*), leftmargin=*]
        \item $\ln\bigl(x\sqrt{x + 1}\bigr)$
        \item $\log_2 \sqrt{\dfrac{x^2 + 1}{x - 1}}$
    \end{enumerate}

    \item Biểu diễn dưới dạng một logarit duy nhất.
    \begin{enumerate}[label=(\alph*), leftmargin=*]
        \item $\frac{1}{2} \ln x - 2\ln(x^2 + 1)$
        \item $\ln(x - 3) + \ln(x + 3) - 2\ln(x^2 - 9)$
    \end{enumerate}
\end{enumerate}

\noindent
\textbf{\textsf{\color{stewartcyan}29--30}}\quad Tìm giá trị chính xác của mỗi biểu thức.

\begin{enumerate}[label=\textbf{\arabic*.}, leftmargin=*, start=29]
    \item 
    \begin{enumerate}[label=(\alph*), itemjoin=\hspace{1.2cm}, leftmargin=*]
        \item $e^{2\ln 5}$
        \item $\log_6 4 + \log_6 54$
        \item $\tan\bigl(\arcsin \frac{4}{5}\bigr)$
    \end{enumerate}

    \item 
    \begin{enumerate}[label=(\alph*), itemjoin=\hspace{1.2cm}, leftmargin=*]
        \item $\ln\dfrac{1}{e^3}$
        \item $\sin(\tan^{-1} 1)$
        \item $10^{-3\log 4}$
    \end{enumerate}
\end{enumerate}

\vspace{0.16cm}
\noindent{\color{stewartcyan}\rule{\linewidth}{0.8pt}}

\end{minipage}%
\hfill
\begin{minipage}[t]{0.485\textwidth}

\noindent
\textbf{\textsf{\color{stewartcyan}31--36}}\quad Giải phương trình tìm $x$. Đưa ra cả giá trị chính xác và giá trị xấp xỉ thập phân, làm tròn đến ba chữ số sau dấu phẩy.

\vspace{0.12cm}
\noindent
\begin{minipage}[t]{0.48\textwidth}
\begin{enumerate}[label=\textbf{\arabic*.}, leftmargin=*, start=31]
    \item $e^{2x} = 3$
    \vspace{0.12cm}
    \item $e^{e^x} = 10$
    \vspace{0.12cm}
    \item $\tan^{-1}(3x^2) = \dfrac{\pi}{4}$
\end{enumerate}
\end{minipage}%
\hfill
\begin{minipage}[t]{0.48\textwidth}
\begin{enumerate}[label=\textbf{\arabic*.}, leftmargin=*, start=32]
    \item $\ln x^2 = 5$
    \vspace{0.12cm}
    \item $\cos^{-1} x = 2$
    \vspace{0.12cm}
    \item $\ln x - 1 = \ln(5 + x) - 4$
\end{enumerate}
\end{minipage}

\vspace{0.16cm}
\noindent{\color{stewartcyan}\rule{\linewidth}{0.8pt}}
\vspace{0.08cm}

\begin{enumerate}[label=\textbf{\arabic*.}, leftmargin=*, start=37]
    \item Tải lượng virus ở một bệnh nhân nhiễm HIV là $52{,}0$ bản sao RNA/mL trước khi bắt đầu điều trị. Tám ngày sau, tải lượng virus giảm còn một nửa so với mức ban đầu.
    \begin{enumerate}[label=(\alph*), leftmargin=*]
        \item Tìm tải lượng virus sau 24 ngày.
        \item Tìm tải lượng virus $V(t)$ còn lại sau $t$ ngày.
        \item Tìm công thức cho hàm ngược của hàm số $V$ và giải thích ý nghĩa của nó.
        \item Sau bao nhiêu ngày thì tải lượng virus sẽ giảm xuống còn $2{,}0$ bản sao RNA/mL?
    \end{enumerate}

    \item Quần thể của một loài nào đó trong môi trường hạn chế với số lượng cá thể ban đầu là $100$ và sức chứa môi trường là $1000$ được mô tả bởi
    \[
    P(t) = \frac{100\,000}{100 + 900e^{-t}}
    \]
    trong đó $t$ được tính bằng năm.

    \begin{tikzpicture}[scale=0.22, baseline=-1mm]
        \draw[stewartcyan, thick] (0,0) rectangle (1.4,1.4);
        \draw[stewartcyan, thin] (0,0.7) -- (1.4,0.7);
        \draw[stewartcyan, thin] (0.7,0) -- (0.7,1.4);
        \draw[stewartcyan, thick] (0.15,0.2) to[out=10, in=230] (0.7,0.7) to[out=50, in=190] (1.25,1.2);
    \end{tikzpicture}
    \begin{enumerate}[label=(\alph*), leftmargin=*]
        \item Vẽ đồ thị của hàm số này và ước lượng cần bao nhiêu thời gian để quần thể đạt đến $900$ cá thể.
        \item Tìm hàm ngược của hàm số này và giải thích ý nghĩa của nó.
        \item Sử dụng hàm ngược để tìm thời gian cần thiết để quần thể đạt đến $900$ cá thể. So sánh với kết quả ở câu (a).
    \end{enumerate}
\end{enumerate}

\end{minipage}